% Stem-Exempt INT8 Quantization — conference manuscript.
% Generated from PAPER_DRAFT.md; every number is checked against results/ by
% scripts/audit_paper_numbers.py (25 claims).
%
% Build:  cd PERCEPTUAL_MP/paper && pdflatex stem_exempt_quantization && pdflatex stem_exempt_quantization
\documentclass[conference]{IEEEtran}

\usepackage{booktabs}
\usepackage{graphicx}
\usepackage{amsmath,amssymb}
\usepackage{textcomp}
\usepackage{url}
\usepackage[hidelinks]{hyperref}

\graphicspath{{../results/phase4/}}

\begin{document}

\title{0.13\,\% of the Weights Carry 99\,\% of the Loss:\\Stem-Exempt INT8 Quantization\\of a Compact Speech Enhancer}

\author{\IEEEauthorblockN{SPARC FYP Team}
\IEEEauthorblockA{\textit{Final Year Project}\\
sparcproject2026@gmail.com}}

\maketitle

\begin{abstract}
Post-training quantization (PTQ) of speech enhancement models is normally applied uniformly, and
mixed-precision variants assign precision by layer \emph{type} or by a second-order sensitivity
proxy. We report a per-layer PTQ study of a 722\,K-parameter Bark-band GTCRN denoiser in which the
allocation is chosen on a speaker-disjoint development split and then evaluated once,
out-of-sample, on the 824-clip VoiceBank-DEMAND test set and on 1\,100 MS-SNSD clips.
We set out to test whether ranking layers by a \emph{perceptual} criterion (PESQ) differs from
ranking them by an MSE-family criterion, and \textbf{it does not}: across the six quantizable
layers the two rankings agree at Spearman $\rho = +0.943$ against SI-SDR and $\rho = +1.000$
against mask-MSE, and the two allocations differ by $0.0003$\,dB. We report this negative result in
full, because it says something useful---for this model class, the cheap MSE-based sensitivity
proxy is already perceptually adequate.
The study's positive finding came from the sensitivity sweep itself. A single \textbf{3.8\,KB stem
convolution---$0.133\,\%$ of the model's weights---carries $98.7\,\%$ of the total INT8
degradation} on the held-out test set. Exempting it produces a model that is statistically
\textbf{indistinguishable from FP32} ($\Delta = 0.0014$\,dB, $p = 0.25$, $n = 824$) at
$3.2\times$ compression, versus a uniform-INT8 loss of $0.109$\,dB
($p = 3.7\mathrm{e}{-51}$). The exemption is significant in \textbf{all nine test conditions} and
replicates on MS-SNSD. Counter-intuitively the mixed model is also \emph{smaller} than uniform INT8
($916.9$ vs $950.7$\,KB), because quantizing every layer causes the runtime to insert QDQ pairs
around 143 additional activation tensors whose graph overhead exceeds the stem's $2.8$\,KB of
weight savings.
We further show that layer-wise quantization errors are \textbf{not additive}, and document two
runtime limits: the band GRU cannot be quantized at all, and INT8 inference is $1.4\times$
\emph{slower} than FP32 on x86.
\end{abstract}

\begin{IEEEkeywords}
post-training quantization, mixed precision, speech enhancement, model compression, ONNX Runtime
\end{IEEEkeywords}

\section{Introduction}

Compact speech enhancement models are deployed under a memory budget, and PTQ is the cheapest way
to meet it---no retraining, no labelled calibration targets. The default is uniform INT8. When that
costs too much quality, the standard remedies are quantization-aware training (expensive) or mixed
precision, where some layers are kept at higher precision.

Mixed precision needs a rule for \emph{which} layers. Two rules dominate. The first assigns by
layer \textbf{type}---for RNN denoisers on microcontrollers, Rusci \emph{et al.}~\cite{rusci2022}
quantize the recurrent layers to INT8 and keep the rest at FP16. The second assigns by a
\textbf{second-order sensitivity proxy}: HAWQ~\cite{dong2019} and HAWQ-V2~\cite{dong2020} rank
layers by top Hessian eigenvalue or average Hessian trace and allocate bits accordingly.

Both proxies are ultimately about reconstruction error. But the objective that matters for speech
enhancement is \emph{perceptual}, and perceptual quality weights frequency bands very
non-uniformly. That gap motivated this study: \textbf{if the perceptual ranking of layers differs
from the MSE ranking, allocating by the perceptual one should win at matched size.}

It does not differ, and we report that plainly (Sec.~\ref{sec:negative}). What the sweep did reveal
is a much sharper structural fact (Sec.~\ref{sec:stem}): the degradation is not spread across
layers at all. It is almost entirely attributable to one tiny layer at the very front of the
network.

\textbf{Contributions.}
\begin{enumerate}
\item \textbf{A negative result, reported in full} (Sec.~\ref{sec:negative}): for this model at
INT8 and 6-layer granularity, perceptual and MSE-family sensitivity rankings coincide. The cheap
proxy is already adequate; practitioners need not build a PESQ-in-the-loop allocator.
\item \textbf{The stem-exemption result} (Sec.~\ref{sec:stem}): $0.133\,\%$ of the weights carry
$98.7\,\%$ of INT8 loss---validated out-of-sample after the allocation was frozen and timestamped.
\item \textbf{A size inversion that contradicts the naive model of quantization}: quantizing
\emph{more} layers produced a \emph{larger} file, because QDQ graph overhead outgrew the weight
savings. This is invisible in any bits-per-weight analysis.
\item \textbf{Condition-level robustness and an explicit cross-corpus caveat}
(Sec.~\ref{sec:conditions}--\ref{sec:crosscorpus}).
\item \textbf{Evidence that layer-wise quantization error is non-additive}
(Sec.~\ref{sec:additive}), which invalidates the additive-cost assumption behind greedy per-layer
allocation search.
\item \textbf{Two hard runtime limits} (Sec.~\ref{sec:runtime}) that bound this problem on ONNX
Runtime.
\end{enumerate}

\section{Related work and positioning}

\textbf{Mixed-precision PTQ for speech enhancement.} The closest prior work is Rusci \emph{et
al.}~\cite{rusci2022}, who deploy an RNN denoiser on a 9-core RISC-V MCU using mixed FP16--INT8
PTQ, reporting $1.4$--$1.7\times$ memory saving at $0.06$ PESQ degradation against $0.3$ for
uniform INT8. Two differences matter. First, their allocation is by \textbf{layer type}---recurrent
layers to INT8---whereas ours is derived from measured per-layer sensitivity. Second, and more
striking, our result runs in the \emph{opposite direction on the same architectural family}: on our
runtime the recurrent layer cannot be quantized at all (Sec.~\ref{sec:runtime}), and the layer that
must be protected is the small input convolution. We read this as evidence that ``which layer to
protect'' is a property of the \emph{runtime and model}, not a transferable rule---which is an
argument for measuring rather than assuming.

\textbf{Hessian-guided allocation.} HAWQ~\cite{dong2019} and HAWQ-V2~\cite{dong2020} rank layers by
curvature and are the standard sensitivity-based allocators. Our Sec.~\ref{sec:negative} result is
adjacent but distinct: we do not compare against a Hessian proxy; we compare a \emph{perceptual}
ranking against \emph{direct output-error} rankings, and find they coincide. Adding a Hessian-trace
ranking as a third column is the single most valuable extension to this work.

\textbf{Speech enhancement compression generally.} Other routes to the same budget include
sign-exponent-only floating point~\cite{hsieh2021}, matrix-product-operator
factorisation~\cite{sun2020}, and structured sparsity for tiny accelerators. These are
complementary to PTQ rather than alternatives to it; our contribution is about \emph{where} the
loss lives under PTQ, not about beating these methods.

\textbf{Objective mismatch.} The framing---that an MSE-optimal quantizer may be wrong for the
objective one actually cares about---was taken from TurboQuant~\cite{zandieh2025}. Its
\emph{algorithm} does not apply here: there is no KV cache, no vector database, no memory
bottleneck of the kind it addresses. We cite it for motivation only, and the perceptual-vs-MSE form
of the hypothesis it suggested was refuted (Sec.~\ref{sec:negative}).

\section{Method}

\textbf{Model.} A Bark-band GTCRN, \texttt{ch96\_h96\_b4}, 722\,209 parameters---a stem
convolution, four gated time--frequency convolution blocks with dilations $1/2/4/8$, a per-band GRU
decoder, and a linear head producing per-band gains. FP32 size $2842.7$\,KB.

\textbf{Quantization.} ONNX Runtime \texttt{quantize\_static}, QDQ format, per-layer control via
\texttt{nodes\_to\_quantize}. Activation calibration uses 64 clips drawn from 11 training speakers.
Per-layer control was verified to produce genuinely distinct graphs (uniform INT8 $=$ 172 QDQ
nodes, head-only $=$ 5, stem-only $=$ 6) with distinct outputs, so the sensitivity study is
mechanically meaningful rather than a no-op.

\textbf{Leakage discipline.} Bit allocation is a model-selection decision, so it is made on a
development set carved from TRAIN \textbf{by speaker}, never on the test set
(Table~\ref{tab:splits}). The split builder hard-fails on speaker overlap, and the chosen
allocation was written to disk \emph{before} any test-set evaluation was run.

\begin{table}[t]
\centering
\caption{Speaker-disjoint splits. The test and external sets were each touched exactly once, after
the allocation was frozen.}
\label{tab:splits}
\begin{tabular}{llrl}
\toprule
split & speakers & clips & used for \\
\midrule
calibration & 11 train spk. & 64 & QDQ activation calibration \\
dev & p254, p256 & 200 & ranking + allocation search \\
test & p232, p257 & 824 & final numbers \\
external & MS-SNSD & 1100 & generalization \\
\bottomrule
\end{tabular}
\end{table}

\textbf{Search space.} The band GRU is unquantizable on this runtime (Sec.~\ref{sec:runtime}) and
INT4 is unavailable for Conv, so the precision axis is per-layer $\in \{$FP32, INT8$\}$ over
\textbf{6} quantizable layers (5 Conv $+$ 1 MatMul)---$2^6 = 64$ configurations, small enough to
enumerate exhaustively rather than search greedily. This matters given the non-additivity in
Sec.~\ref{sec:additive}.

\section{Experimental setup}

VoiceBank-DEMAND 16\,kHz; 824-clip test set with official condition labels (5 noise types
$\times$ 4 input SNRs). MS-SNSD provides 1\,100 external clips. Metrics: SI-SDR, STOI, ESTOI, LSD,
and PESQ. Significance throughout is paired Wilcoxon signed-rank with $95\,\%$ bootstrap CIs on
per-clip differences.

\section{Results}

\subsection{Perceptual and MSE sensitivity rankings coincide (negative result)}
\label{sec:negative}

\begin{table}[t]
\centering
\caption{Rank correlation between sensitivity orderings over the six quantizable layers. The
perceptual and MSE-family rankings agree almost perfectly.}
\label{tab:rankings}
\begin{tabular}{lr}
\toprule
ranking comparison & Spearman $\rho$ \\
\midrule
PESQ vs SI-SDR & $+0.943$ \\
PESQ vs mask-MSE & $\mathbf{+1.000}$ \\
\bottomrule
\end{tabular}
\end{table}

\texttt{stem} ranks most sensitive under \emph{every} metric; \texttt{head} ranks least sensitive
under \emph{every} metric (Table~\ref{tab:rankings}). The two induced allocations differ at exactly
one position and produce a $\Delta$SI-SDR of $+0.0003$\,dB---indistinguishable.

\textbf{Conclusion:} at INT8 and this granularity, MSE-guided allocation is already perceptually
adequate for this model. We make no perceptual-advantage claim.

\subsection{The stem exemption (main result, out-of-sample)}
\label{sec:stem}

\begin{table}[t]
\centering
\caption{Post-training quantization on the VoiceBank-DEMAND test set (824 clips, held out from the
dev split used to choose the allocation). Exempting the 3.8\,KB stem convolution recovers
essentially all of the quantization loss and yields a model statistically indistinguishable from
FP32.}
\label{tab:main}
\begin{tabular}{lrrr}
\toprule
Method & Size (KB) & SI-SDR (dB) & vs FP32 \\
\midrule
FP32 (no quantization) & $2842.7$ & $18.377$ & --- \\
\textbf{Mixed: stem FP32 (ours)} & $\mathbf{916.9}$ & $\mathbf{18.375}$ & $p = 0.25$ (n.s.) \\
Uniform INT8 & $950.7$ & $18.267$ & $p = 3.7\mathrm{e}{-51}$ \\
\bottomrule
\end{tabular}
\end{table}

From Table~\ref{tab:main}:
\begin{itemize}
\item Stem exemption vs uniform: $+0.108$\,dB SI-SDR, $95\,\%$ CI $[+0.091, +0.126]$,
$p = 9.2\mathrm{e}{-49}$.
\item Mixed vs FP32: $\Delta = 0.0014$\,dB, CI $[-0.0012, +0.0041]$, \textbf{not significant}.
\item Total quantization cost (FP32 $-$ uniform) is $0.109$\,dB; the stem accounts for
$\mathbf{98.7\,\%}$ of it.
\item The stem is $\mathbf{3\,840}$ bytes $= 0.133\,\%$ of model weights.
\end{itemize}

\textbf{The mixed model is also smaller than uniform INT8.} Verified cause: requesting quantization
of every layer makes ONNX Runtime wrap ${\sim}143$ additional activation tensors in QDQ pairs
(\textbf{172 vs 29} QDQ nodes), and that graph overhead exceeds the stem's $2.8$\,KB of weight
savings. Quantizing the stem \emph{alone} even grows the file above FP32 ($2851.9$ vs
$2842.7$\,KB). Any analysis that counts only bits-per-weight will miss this entirely.

\textbf{Fidelity to FP32}---mean absolute per-clip deviation, the property PTQ is actually supposed
to deliver---favours the mixed allocation in \textbf{8/8 metric--corpus cells, by $7$--$20\times$}
(Table~\ref{tab:fidelity}).

\begin{table}[t]
\centering
\caption{Mean absolute per-clip deviation from FP32 (lower is more faithful). All eight cells
favour the mixed allocation.}
\label{tab:fidelity}
\begin{tabular}{llrrr}
\toprule
corpus & metric & mixed & uniform & adv. \\
\midrule
VB-824 & $|\Delta$SI-SDR$|$ & $\mathbf{0.0202}$ & $0.1458$ & $7.2\times$ \\
VB-824 & $|\Delta$STOI$|$ & $\mathbf{0.00025}$ & $0.00186$ & $7.4\times$ \\
VB-824 & $|\Delta$ESTOI$|$ & $\mathbf{0.00055}$ & $0.00481$ & $8.7\times$ \\
VB-824 & $|\Delta$LSD$|$ & $\mathbf{0.0088}$ & $0.1708$ & $19.5\times$ \\
MS-SNSD & $|\Delta$SI-SDR$|$ & $\mathbf{0.0328}$ & $0.2832$ & $8.6\times$ \\
MS-SNSD & $|\Delta$STOI$|$ & $\mathbf{0.00051}$ & $0.00710$ & $13.9\times$ \\
MS-SNSD & $|\Delta$ESTOI$|$ & $\mathbf{0.00079}$ & $0.00966$ & $12.2\times$ \\
MS-SNSD & $|\Delta$LSD$|$ & $\mathbf{0.0209}$ & $0.2915$ & $13.9\times$ \\
\bottomrule
\end{tabular}
\end{table}

\subsection{Condition robustness}
\label{sec:conditions}

\begin{table}[t]
\centering
\caption{Benefit of exempting the stem, by noise type and input SNR. The effect is statistically
significant in every one of the nine conditions.}
\label{tab:conditions}
\begin{tabular}{lrrl}
\toprule
Condition & $n$ & $\Delta$SI-SDR (dB) & 95\,\% CI \\
\midrule
bus & 164 & $+0.087$ & $[+0.051, +0.130]$ \\
cafe & 164 & $+0.156$ & $[+0.124, +0.189]$ \\
living & 166 & $+0.098$ & $[+0.064, +0.134]$ \\
office & 164 & $+0.024$ & $[+0.011, +0.038]$ \\
psquare & 166 & $+0.175$ & $[+0.126, +0.235]$ \\
\midrule
2.5\,dB & 207 & $+0.125$ & $[+0.088, +0.165]$ \\
7.5\,dB & 206 & $+0.095$ & $[+0.065, +0.128]$ \\
12.5\,dB & 207 & $+0.104$ & $[+0.076, +0.137]$ \\
17.5\,dB & 204 & $+0.108$ & $[+0.075, +0.151]$ \\
\bottomrule
\end{tabular}
\end{table}

The effect is \textbf{flat across SNR} (Table~\ref{tab:conditions})---so it is not an SNR
artifact---and largest in hard non-stationary noise (psquare, cafe), smallest in easy stationary
office noise. That ordering is what one would predict if the stem's early feature extraction
matters most when spectral structure is complex, which is a satisfying but \emph{post-hoc}
explanation; we offer it as interpretation, not evidence.

\subsection{Cross-corpus caveat---stated explicitly}
\label{sec:crosscorpus}

On MS-SNSD the stem-exemption SI-SDR benefit replicates ($+0.039$\,dB, $p = 3.0\mathrm{e}{-12}$)
and mixed remains indistinguishable from FP32 on SI-SDR ($p = 0.17$). \textbf{However, on MS-SNSD
uniform INT8 scores better than mixed on STOI ($-0.0060$), ESTOI ($-0.0067$) and LSD ($+0.233$).}

This is not evidence that uniform is the better allocation. Uniform INT8 also beats \textbf{FP32
itself} on MS-SNSD STOI ($0.8212$ vs $0.8154$)---a quantizer cannot improve on the model it is
approximating. Its larger quantization noise acts as a mild dither that flatters intelligibility
metrics under distribution shift while degrading waveform fidelity. The mixed model instead tracks
FP32 closely on every metric (Table~\ref{tab:fidelity}).

\textbf{The claim we make is therefore fidelity to FP32, not ``better enhancement on every metric
everywhere.''} We state this in the abstract as well as here, because the temptation to overclaim
is exactly what this result guards against.

\subsection{Quantization error is not additive}
\label{sec:additive}

\begin{table}[t]
\centering
\caption{Layer-wise degradations do not sum. The interaction term's CI excludes zero on both
corpora.}
\label{tab:additive}
\begin{tabular}{lrrrr}
\toprule
corpus & stem & 5 blocks & combined & ratio \\
\midrule
dev (200) & $-0.0061$ & $-0.0023$ & $-0.0996$ & $\mathbf{11.95\times}$ \\
MS-SNSD (1100) & $+0.0260$ & $+0.0010$ & $-0.0377$ & $1.40\times$ \\
\bottomrule
\end{tabular}
\end{table}

Error \textbf{compounds} through the network rather than accumulating independently
(Table~\ref{tab:additive}). The magnitude is strongly corpus-dependent, so we report the
\emph{existence} of the interaction as robust and its \emph{size} as corpus-specific.

Practical consequence: greedy per-layer allocation search assumes additive costs and is therefore
unsound here. Exhaustive enumeration over the 64 configurations was affordable and is what we did.

\subsection{Runtime limits (findings, not bugs)}
\label{sec:runtime}

\begin{itemize}
\item \textbf{The band GRU cannot be quantized} by ONNX Runtime 1.20.1. Requesting it yields
\textbf{0 QDQ nodes}, no size change ($2852.9$ vs $2842.7$\,KB) and no speedup ($0.995\times$). The
recurrent decoder is incompressible on this runtime and is pinned to FP32 in every allocation. Note
the direct contrast with~\cite{rusci2022}, whose method depends on quantizing exactly this layer
type.
\item \textbf{INT4 is unavailable for Conv} via \texttt{quantize\_static}. The runtime's 4-bit path
covers MatMul only, and this graph has one tiny MatMul, so INT4 cannot reach the weight-heavy Conv
layers.
\item \textbf{INT8 is ${\sim}1.4\times$ slower than FP32} on x86, re-measured on a verified-idle
machine ($6\,\%$ CPU load, AMD Ryzen~5 5600H, same runtime version): $1.39\times$ / $1.43\times$ /
$1.46\times$ at 1 / 2 / 4 threads, with absolute RTF $0.084$ (FP32) vs $0.117$ (INT8)
single-threaded. \textbf{On this hardware the benefit of quantization is memory, not speed}---and
the disadvantage grows slightly with thread count, consistent with QDQ overhead parallelising worse
than the dense kernels it replaces.
\end{itemize}

\section{Limitations}

The negative result in Sec.~\ref{sec:negative} is scoped to this checkpoint, at INT8, over six
layers. With a finer precision axis (INT4 available, or per-channel), or a model whose layers are
less dominated by a single one, the perceptual and MSE rankings could plausibly diverge. We did not
show that they never diverge; we showed they do not here, and explained why the search space was
forced to be this coarse (Sec.~\ref{sec:runtime}).

One checkpoint, one architecture. The stem-exemption magnitude is not claimed to transfer. The
transferable claim is methodological: measure per-layer sensitivity, because the answer is
runtime- and model-specific and contradicted a published rule for the same architecture family.

No Hessian baseline. Without it, Sec.~\ref{sec:negative} compares a perceptual ranking against
direct output-error rankings, not against the standard curvature proxy.

No on-device measurement. All timings are x86 with thread pinning as an edge proxy. The contrast
with~\cite{rusci2022} is therefore across both method \emph{and} hardware, and we cannot separate
those two factors.

\section{Conclusion}

We tested whether perceptually-guided bit allocation beats MSE-guided allocation for a compact
speech enhancer, and found that it does not---the rankings coincide almost perfectly. Reporting
that negative result is worthwhile on its own: it licenses practitioners to keep using the cheap
proxy.

The sweep that produced it revealed something more useful. Quantization loss in this model is not
distributed across layers; $\mathbf{0.133\,\%}$ \textbf{of the weights carry} $\mathbf{98.7\,\%}$
\textbf{of it}. Protecting that one 3.8\,KB layer yields an FP32-indistinguishable model at
$3.2\times$ compression, significant in every test condition, replicated on a second corpus---and,
because of QDQ graph overhead, \emph{smaller} than quantizing everything. Combined with the
non-additivity of layer-wise errors and the fact that a published mixed-precision rule for this
architecture family is inverted on our runtime, the practical lesson is consistent: for PTQ of
small speech models, per-layer sensitivity is worth measuring rather than assuming.

\begin{thebibliography}{99}

\bibitem{rusci2022}
M.~Rusci, M.~Fariselli, M.~Croome, F.~Paci, and E.~Flamand, ``Accelerating RNN-based speech
enhancement on a multi-core MCU with mixed FP16-INT8 post-training quantization,'' in \emph{ITEM
Workshop, ECML-PKDD}, 2022.

\bibitem{dong2019}
Z.~Dong, Z.~Yao, A.~Gholami, M.~W.~Mahoney, and K.~Keutzer, ``HAWQ: Hessian aware quantization of
neural networks with mixed-precision,'' in \emph{Proc.\ ICCV}, 2019.

\bibitem{dong2020}
Z.~Dong, Z.~Yao, D.~Arfeen, A.~Gholami, M.~W.~Mahoney, and K.~Keutzer, ``HAWQ-V2: Hessian aware
trace-weighted quantization of neural networks,'' in \emph{Proc.\ NeurIPS}, 2020.

\bibitem{zandieh2025}
A.~Zandieh, M.~Daliri, M.~Hadian, and V.~Mirrokni, ``TurboQuant: Online vector quantization with
near-optimal distortion rate,'' \emph{arXiv:2504.19874}, 2025.

\bibitem{hsieh2021}
Y.-T.~Hsieh \emph{et al.}, ``SEOFP-NET: Compression and acceleration of deep neural networks for
speech enhancement using sign-exponent-only floating-points,'' \emph{arXiv:2111.04436}, 2021.

\bibitem{sun2020}
X.~Sun \emph{et al.}, ``A model compression method with matrix product operators for speech
enhancement,'' \emph{arXiv:2010.04950}, 2020.

\end{thebibliography}

\end{document}
